% (User input window)

in proposition \ref{prop:pom-diagonal-rate-keep-resample-channel} In addition to the hold-resampling description, the optimal kernel can also be equivalently rewritten as a rank-one perturbation superposed by state-dependent laziness and global refresh.
This representation enables characteristic polynomials with finite dimensions \(\zeta\) structured decomposition of the determinant becomes a direct corollary, and in \(\delta\downarrow 0\), a class of auditable critical spectral operators is derived.

\begin{proposition}[Rank one refresh core representation]\label{prop:pom-diagonal-rate-rank-one-refresh-decomposition}
inheritance proposition \ref{prop:pom-diagonal-rate-keep-resample-channel} The sign of
\(
A:=\sum_{z\in X}u_z
\)
and define the refresh rate
\[
r_\delta(x):=\frac{A}{A+\kappa u_x}\in(0,1),\qquad x\in X.
\]
to everything \(x,y\in X\) has a strict identity
\[
K_\delta(y\mid x)
=r_\delta(x)\,\pi_\delta(y)+\bigl(1-r_\delta(x)\bigr)\mathbf 1_{\{y=x\}}.
\]
So in matrix form
\[
K_\delta=\mathrm{diag}(1-r_\delta)+r_\delta\,\pi_\delta^{\mathsf T},
\]
in \(r_\delta,\pi_\delta\) are treated as column vectors.
\end{proposition}
\begin{proof}
by proposition \ref{prop:pom-diagonal-rate-keep-resample-channel} proof, for any \(y\ne x\) have
\[
K_\delta(y\mid x)=\frac{u_y}{A+\kappa u_x}
=\frac{A}{A+\kappa u_x}\cdot \frac{u_y}{A}
=r_\delta(x)\,\pi_\delta(y).
\]
The diagonal term is given by
\[
K_\delta(x\mid x)=\frac{(1+\kappa)u_x}{A+\kappa u_x}
=1-\frac{A-u_x}{A+\kappa u_x}
=1-r_\delta(x)\bigl(1-\pi_\delta(x)\bigr)
=r_\delta(x)\pi_\delta(x)+\bigl(1-r_\delta(x)\bigr)
\]
Obtain the representation shown.
The matrix form is a direct result of row-by-row rewriting. Certification completed.
\end{proof}

\begin{theorem}[Characteristic polynomials and finite dimensions \texorpdfstring{$\zeta$}{zeta} One-dimensional collapse of determinant]\label{thm:pom-diagonal-rate-rank-one-refresh-determinant-collapse}
in proposition \ref{prop:pom-diagonal-rate-rank-one-refresh-decomposition} Under the setting of
\(
D_\delta:=\mathrm{diag}(1-r_\delta)
\)。
Then for any \(z\in\CC\) make \(I-zD_\delta\) reversible, with determinant decomposition
\[
\det(I-zK_\delta)=\det(I-zD_\delta)\Bigl(1-z\,\pi_\delta^{\mathsf T}(I-zD_\delta)^{-1}r_\delta\Bigr).
\]
Equivalently,
\[
\det(I-zK_\delta)
=\Bigl(\prod_{x\in X}\bigl(1-z(1-r_\delta(x))\bigr)\Bigr)
\Bigl(1-z\sum_{x\in X}\frac{\pi_\delta(x)r_\delta(x)}{1-z(1-r_\delta(x))}\Bigr).
\]
Therefore after the diagonal parts are stripped away, the remaining spectrum is determined by a scalar secular equation.
\end{theorem}
\begin{proof}
by proposition \ref{prop:pom-diagonal-rate-rank-one-refresh-decomposition},have \(K_\delta=D_\delta+r_\delta\pi_\delta^{\mathsf T}\)。
right
\(
I-zK_\delta=(I-zD_\delta)-z r_\delta\pi_\delta^{\mathsf T}
\)
Apply the rank-determinant lemma
\(
\det(D+uv^{\mathsf T})=\det(D)\bigl(1+v^{\mathsf T}D^{-1}u\bigr)
\)
That’s the first formula.
because \(I-zD_\delta\) is a diagonal matrix,\(\det(I-zD_\delta)=\prod_x(1-z(1-r_\delta(x)))\),and
\[
\pi_\delta^{\mathsf T}(I-zD_\delta)^{-1}r_\delta
=\sum_{x\in X}\pi_\delta(x)\frac{r_\delta(x)}{1-z(1-r_\delta(x))},
\]
Thus the second formula is obtained. Certification completed.
\end{proof}

\begin{corollary}[Thermal Traces and Finite Dimensions \texorpdfstring{$\zeta$}{zeta} Fingerprint: Rational and Extreme Reality]\label{cor:pom-diagonal-rate-rank-one-refresh-trace-zeta-fingerprint}
make \(n:=|X|\ge 2\), and note \(K_\delta\) eigenvalues ​​(including multiplicity) are
\(
1=\lambda_1(K_\delta),\lambda_2(K_\delta),\dots,\lambda_n(K_\delta)
\)。
Define finite dimensions \(\zeta\) determinant
\[
\zeta_\delta(z):=\det(I-zK_\delta)^{-1}.
\]
but \(\zeta_\delta\) is a rational function, and its pole set is
\[
\mathrm{Pole}(\zeta_\delta)=\Bigl\{\frac{1}{\lambda_i(K_\delta)}:i=1,\dots,n\Bigr\},
\]
and all poles fall on the real axis; more precisely, by \(\mathrm{Spec}(K_\delta)\subseteq[0,1]\)(inference \ref{cor:pom-maxent-markov-no-oscillation-endpoints})have to
\(
\mathrm{Pole}(\zeta_\delta)\subseteq [1,\infty]
\)。
\par
In addition, in \(z=0\) The identity of the neighborhood trace generating function of
\[
\log \zeta_\delta(z)=\sum_{m\ge 1}\frac{\tr(K_\delta^m)}{m}\,z^m,
\qquad
\tr(K_\delta^m)=\sum_{i=1}^{n}\lambda_i(K_\delta)^m\quad(m\ge 1).
\]
therefore \((\tr(K_\delta^m))_{m\ge 1}\) can be determined by the theorem \ref{thm:pom-diagonal-rate-rank-one-refresh-determinant-collapse} The closed form of the one-dimensional determinant of is recovered item by item through coefficient extraction,
and as \((r_\delta,\pi_\delta)\) is a completely symmetric fingerprint (invariant to label permutation).
\end{corollary}
\begin{proof}
From finite-dimensional linear algebra,\(\det(I-zK_\delta)=\prod_{i=1}^{n}(1-z\lambda_i(K_\delta))\) is the number of times \(n\) polynomial,
Therefore \(\zeta_\delta\) is a rational function with poles \(\{1/\lambda_i\}\)。
Spectral reality is symmetrized by similarity \(S_\delta=D^{1/2}K_\delta D^{-1/2}\)(theorem \ref{thm:pom-diagonal-rate-critical-line-laplacian-first-order})and \(S_\delta\) symmetry is given;
Spectral non-negative AND \(\le 1\) has been deduced \ref{cor:pom-maxent-markov-no-oscillation-endpoints} Proved.
\par
Finally, from the power series identity \(-\log(1-t)=\sum_{m\ge 1}t^m/m\) and the spectral representation of the trace
\(
\tr(K_\delta^m)=\sum_i \lambda_i(K_\delta)^m
\)
Directly derive the trace generating function formula.
Will \(\zeta_\delta=\det(I-zK_\delta)^{-1}\) and theorem \ref{thm:pom-diagonal-rate-rank-one-refresh-determinant-collapse} closed combination, you can press \(z\) of the Taylor coefficient item by item. \(\tr(K_\delta^m)\), and since the expression is in \(x\) is completely symmetrical and invariant. Certification completed.
\end{proof}

\begin{corollary}[Closed form of determinant to remove trivial eigenvalues]\label{cor:pom-diagonal-rate-rank-one-refresh-detprime}
make \(n:=|X|\ge 2\), and order \(1=\lambda_1(K_\delta),\lambda_2(K_\delta),\dots,\lambda_n(K_\delta)\) for \(K_\delta\) eigenvalues ​​(including multiplicity).
definition
\[
\det{}'(I-K_\delta):=\prod_{i=2}^{n}\bigl(1-\lambda_i(K_\delta)\bigr)
=\lim_{z\to 1}\frac{\det(I-zK_\delta)}{1-z}.
\]
Then there is a strict identity
\[
\det{}'(I-K_\delta)
=\Bigl(\prod_{x\in X}r_\delta(x)\Bigr)\Bigl(\sum_{x\in X}\frac{\pi_\delta(x)}{r_\delta(x)}\Bigr).
\]
And by the marginal self-consistent relationship \(w(x)=u_x(A+\kappa u_x)/Z\) roll out
\[
\sum_{x\in X}\frac{\pi_\delta(x)}{r_\delta(x)}=\frac{Z}{A^2}.
\]
If we take the theorem \ref{thm:pom-diagonal-rate-scalar-collapse} normalized scaling of (thus \(Z=1\)), then further simplifies to
\[
\det{}'(I-K_\delta)=\frac{\prod_{x\in X}r_\delta(x)}{A^2}.
\]
\end{corollary}
\begin{proof}
Depend on \(\det(I-zK_\delta)=\prod_{i=1}^{n}(1-z\lambda_i(K_\delta))\) and \(\lambda_1(K_\delta)=1\),have to
\(
\det{}'(I-K_\delta)=\lim_{z\to 1}\det(I-zK_\delta)/(1-z)
\)。
\par
will theorem \ref{thm:pom-diagonal-rate-rank-one-refresh-determinant-collapse} Substitute and record
\[
g(z):=\det(I-zD_\delta)=\prod_{x\in X}\bigl(1-z(1-r_\delta(x))\bigr),
\qquad
s(z):=\pi_\delta^{\mathsf T}(I-zD_\delta)^{-1}r_\delta.
\]
then in \(z\to 1\) sometimes \(g(z)\to g(1)=\prod_x r_\delta(x)\),and
\(
1-zs(z)\to 1-s(1)=0
\)(because \(s(1)=\sum_x\pi_\delta(x)=1\)）。
therefore
\[
\det{}'(I-K_\delta)=g(1)\cdot\lim_{z\to 1}\frac{1-zs(z)}{1-z}
=g(1)\cdot\bigl(s(1)+s'(1)\bigr).
\]
Notice
\[
s'(z)=\sum_{x\in X}\pi_\delta(x)\cdot
\frac{r_\delta(x)\bigl(1-r_\delta(x)\bigr)}{\bigl(1-z(1-r_\delta(x))\bigr)^2},
\]
thereby
\(
s'(1)=\sum_x\pi_\delta(x)\frac{1-r_\delta(x)}{r_\delta(x)}
=\sum_x\frac{\pi_\delta(x)}{r_\delta(x)}-1
\)。
Get it on behalf of others
\(
\det{}'(I-K_\delta)
=(\prod_x r_\delta(x))\sum_x \pi_\delta(x)/r_\delta(x)
\)。
\par
Finally, by \(\pi_\delta(x)=u_x/A\)、\(r_\delta(x)=A/(A+\kappa u_x)\) have to
\[
\frac{\pi_\delta(x)}{r_\delta(x)}
=\frac{u_x}{A}\cdot \frac{A+\kappa u_x}{A}
=\frac{u_x(A+\kappa u_x)}{A^2}
=\frac{Zw(x)}{A^2},
\]
Sum and use \(\sum_x w(x)=1\) have to \(\sum_x\pi_\delta(x)/r_\delta(x)=Z/A^2\)。
exist \(Z=1\) is normalized and the conclusion is as shown. Certification completed.
\end{proof}

\begin{theorem}[Critical line Laplacian vs. \texorpdfstring{$\delta\downarrow0$}{delta->0} first-order notation]\label{thm:pom-diagonal-rate-critical-line-laplacian-first-order}
Take the full support distribution \(w\in\Delta(X)\),make \(K_\delta\) for \(P^\star_\delta\) induced optimal kernel (proposition \ref{prop:pom-diagonal-rate-keep-resample-channel}）。
remember \(D:=\mathrm{diag}(w)\), and define the symmetrization operator
\[
S_\delta:=D^{1/2}K_\delta D^{-1/2}=D^{-1/2}P^\star_\delta D^{-1/2},
\]
but \(S_\delta\) and \(K_\delta\) are similar, the spectra are the same and \(S_\delta\) is a real symmetric matrix.
\par
make
\[
A:=A_{1/2}(w)=\sum_{x\in X}\sqrt{w(x)},\qquad
C:=C_{1/2}(w)=A^2-1\ (>0),
\]
And order \(\sqrt w\in\RR^X\) means that the component is \(\sqrt{w(x)}\) column vector.
then when \(\delta\downarrow 0\) when there is a symmetric positive semidefinite matrix
\[
L_w:=\frac{A\,\mathrm{diag}(w^{-1/2})-\mathbf 1\mathbf 1^{\mathsf T}}{C}\succeq 0,
\qquad
L_w\sqrt w=0,
\]
make
\[
S_\delta=I-\delta L_w+o(\delta)\qquad(\delta\downarrow 0).
\]
If you remember \(0=\mu_1<\mu_2\le\cdots\le\mu_n\) for \(L_w\), then
\[
\lambda_i(K_\delta)=1-\delta\mu_i+o(\delta)\qquad(i=1,\dots,n),
\]
in \(\lambda_1(K_\delta)\equiv 1\)。
\end{theorem}
\begin{proof}
Depend on \(w(x)K_\delta(y\mid x)=P^\star_\delta(x,y)=P^\star_\delta(y,x)=w(y)K_\delta(x\mid y)\) have to \(S_\delta\) symmetry.
\par
By theorem \ref{thm:pom-diagonal-rate-rare-mismatch-hellinger-universality},when \(\delta\downarrow 0\) for any \(x\ne y\),
\[
K_\delta(y\mid x)=\frac{P^\star_\delta(x,y)}{w(x)}
=\delta\,\frac{\sqrt{w(y)}}{\sqrt{w(x)}\,C}+o(\delta),
\]
and
\[
K_\delta(x\mid x)=\frac{P^\star_\delta(x,x)}{w(x)}
=1-\delta\,\frac{A-\sqrt{w(x)}}{\sqrt{w(x)}\,C}+o(\delta)
=1-\delta\,\frac{A/\sqrt{w(x)}-1}{C}+o(\delta).
\]
thus to \(x\ne y\),
\[
S_\delta(x,y)
=\frac{\sqrt{w(x)}}{\sqrt{w(y)}}K_\delta(y\mid x)
=\delta\,\frac{1}{C}+o(\delta),
\qquad
S_\delta(x,x)=K_\delta(x\mid x)=1-\delta\,\frac{A/\sqrt{w(x)}-1}{C}+o(\delta).
\]
Expanding and merging the above elements element by element, we get \(S_\delta=I-\delta L_w+o(\delta)\),in \(L_w\)'s diagonal element is \(\frac{A/\sqrt{w(x)}-1}{C}\), the off-diagonal element is \(-\frac{1}{C}\), which is the matrix expression shown.
\par
as evidence \(L_w\succeq 0\),remember \(t_x:=A/\sqrt{w(x)}\) and write
\(
CL_w=\mathrm{diag}(t)-\mathbf 1\mathbf 1^{\mathsf T}
\)。
to any \(v\in\RR^X\), obtained from the weighted Cauchy--Schwarz inequality
\[
\biggl(\sum_{x\in X}v_x\biggr)^2
=\biggl(\sum_{x\in X}\sqrt{\frac{1}{t_x}}\cdot \sqrt{t_x}\,v_x\biggr)^2
\le \biggl(\sum_{x\in X}\frac{1}{t_x}\biggr)\biggl(\sum_{x\in X}t_xv_x^2\biggr).
\]
and
\(
\sum_x\frac{1}{t_x}=\sum_x\frac{\sqrt{w(x)}}{A}=1
\), so
\(
v^{\mathsf T}(\mathrm{diag}(t)-\mathbf 1\mathbf 1^{\mathsf T})v
=\sum_x t_xv_x^2-(\sum_x v_x)^2\ge 0
\),Right now \(L_w\succeq 0\)。
The equal sign holds true if and only if \(\sqrt{t_x}v_x\) and \(\sqrt{1/t_x}\) is proportional to \(v_x\propto 1/t_x=\sqrt{w(x)}/A\),thereby \(L_w\sqrt w=0\)。
\par
Finally, the operator norm is continuous and \(S_\delta=I-\delta L_w+o(\delta)\), available
\(
\lambda_i(S_\delta)=1-\delta\mu_i+o(\delta)
\)；
And because \(S_\delta\) and \(K_\delta\) are similar, the two spectra are the same, and the spectral expansion shown is obtained. Certification completed.
\end{proof}

On top of the first-order spectral law of the theorem \ref{thm:pom-diagonal-rate-critical-line-laplacian-first-order}, the optimal kernel of the small distortion endpoint can be regarded as the Euler discretization of a family of reversible continuous time chains.

\begin{corollary}[Critical continuous time generator and complete graph conductance]\label{cor:pom-diagonal-rate-critical-continuous-time-generator}
Follow the setting and notation of the theorem \ref{thm:pom-diagonal-rate-critical-line-laplacian-first-order}.
define matrix
\[
Q_w:=\lim_{\delta\downarrow 0}\frac{K_\delta-I}{\delta},
\]
Then the limit exists in the sense of matrix elements and holds for any \(x\ne y\)
\[
(Q_w)_{xy}=\frac{1}{C}\sqrt{\frac{w(y)}{w(x)}},
\qquad
(Q_w)_{xx}=-\sum_{y\ne x}(Q_w)_{xy}
=-\frac{A-\sqrt{w(x)}}{C\sqrt{w(x)}}.
\]
Thus \(K_\delta=I+\delta Q_w+o(\delta)\)(\(\delta\downarrow 0\)).
And \(Q_w\) is a reversible continuous time generator with \(w\) as a stationary distribution:
\[
\sum_{y\in X}(Q_w)_{xy}=0,\qquad
w(x)(Q_w)_{xy}=w(y)(Q_w)_{yx}\quad(\forall x,y).
\]
If we define conductance
\[
c_{xy}:=w(x)(Q_w)_{xy}=\frac{1}{C}\sqrt{w(x)w(y)}=c_{yx},
\]
Then the invertible representation of \(Q_w\) on the complete graph is \((Q_w)_{xy}=c_{xy}/w(x)\).
\end{corollary}
\begin{proof}
From the proof of the theorem \ref{thm:pom-diagonal-rate-critical-line-laplacian-first-order} we can get the element-wise expansion:
when \(x\ne y\)
\(
K_\delta(y\mid x)=\delta\,\frac{\sqrt{w(y)}}{\sqrt{w(x)}\,C}+o(\delta)
\)，
and
\(
K_\delta(x\mid x)=1-\delta\,\frac{A-\sqrt{w(x)}}{C\sqrt{w(x)}}+o(\delta)
\)。
Therefore \(\lim_{\delta\downarrow 0}(K_\delta-I)/\delta\) exists element-wise and gives the matrix elements of \(Q_w\) shown.
From \(\sum_y K_\delta(y\mid x)=1\) we get \(\sum_y (Q_w)_{xy}=0\);
The detailed balance is obtained by \(P^\star_\delta(x,y)=w(x)K_\delta(y\mid x)=w(y)K_\delta(x\mid y)\).
The expression for conductance is a direct corollary of the definition. Certification completed.
\end{proof}

\begin{corollary}[Conductance kernel multiplicative closure under tensor product]\label{cor:pom-diagonal-rate-critical-conductance-tensor-multiplicativity}
Take the two complete support distributions \(w\in\Delta(X)\), \(v\in\Delta(Y)\), and let \(\omega:=w\otimes v\in\Delta(X\times Y)\).
remember
\[
A(w):=\sum_{x\in X}\sqrt{w(x)},\quad C(w):=A(w)^2-1,
\qquad
A(v):=\sum_{y\in Y}\sqrt{v(y)},\quad C(v):=A(v)^2-1,
\]
And the corresponding definitions of \(A(\omega),C(\omega)\).
but
\[
A(\omega)=A(w)A(v),\qquad
1+C(\omega)=(1+C(w))(1+C(v)),
\qquad
C(\omega)=C(w)+C(v)+C(w)C(v).
\]
Let \(Q_w,Q_v,Q_\omega\) be the critical continuous time generator corresponding to the inference \ref{cor:pom-diagonal-rate-critical-continuous-time-generator} respectively, and record its conductance as
\[
c^w_{x_1x_2}:=w(x_1)(Q_w)_{x_1x_2},\qquad
c^v_{y_1y_2}:=v(y_1)(Q_v)_{y_1y_2},\qquad
c^\omega_{(x_1,y_1),(x_2,y_2)}:=\omega(x_1,y_1)(Q_\omega)_{(x_1,y_1),(x_2,y_2)}.
\]
Then there is a strict identity for any \((x_1,y_1)\ne(x_2,y_2)\)
\[
c^\omega_{(x_1,y_1),(x_2,y_2)}
=\frac{1}{C(\omega)}\sqrt{w(x_1)w(x_2)}\sqrt{v(y_1)v(y_2)}
=\frac{C(w)C(v)}{C(\omega)}\,c^w_{x_1x_2}\,c^v_{y_1y_2}.
\]
\end{corollary}
\begin{proof}
From \(A(\omega)=A(w)A(v)\), we immediately get \(1+C(\omega)=A(\omega)^2=(1+C(w))(1+C(v))\), which can be expanded to get \(C(\omega)=C(w)+C(v)+C(w)C(v)\).
The conductance identity is derived from the closed form of \ref{cor:pom-diagonal-rate-critical-continuous-time-generator}
\(
c_{xy}=\frac{1}{C}\sqrt{w(x)w(y)}
\)
Substituting \(w,v,\omega\) respectively and sorting them out. Certification completed.
\end{proof}

\begin{theorem}[Entropy maximum reversible generator under fixed average jump rate constraint]\label{thm:pom-diagonal-rate-critical-continuous-time-generator-maxent}
Let \(X\) be a finite set, and take the fully supported distribution \(w\in\Delta(X)\).
Let \(\mathcal Q_{\mathrm{rev}}(w)\) be the set of all reversible continuous time generators with \(w\) as a stationary distribution.
Given normalization constraints
\[
\sum_{x\in X}w(x)\sum_{y\ne x}Q_{xy}=1,
\]
Maximize the functional on \(\mathcal Q_{\mathrm{rev}}(w)\)
\[
\mathsf H(Q):=-\sum_{x\in X}w(x)\sum_{y\ne x}Q_{xy}\log Q_{xy}
\]
The unique solution of is exactly \(Q_w\) of the corollary \ref{cor:pom-diagonal-rate-critical-continuous-time-generator}.
\end{theorem}
\begin{proof}
For any \(Q\in\mathcal Q_{\mathrm{rev}}(w)\) let the symmetric conductance \(c_{xy}:=w(x)Q_{xy}=w(y)Q_{yx}\)(\(x\ne y\)).
Then the constraint becomes
\(
\sum_{x\ne y}c_{xy}=2\sum_{x<y}c_{xy}=1
\)。
and
\[
\mathsf H(Q)
=-\sum_{x\ne y}c_{xy}\log\frac{c_{xy}}{w(x)}
=-\sum_{x<y}c_{xy}\Bigl(\log\frac{c_{xy}}{w(x)}+\log\frac{c_{xy}}{w(y)}\Bigr)
=-\sum_{x<y}2c_{xy}\log c_{xy}+\sum_{x<y}c_{xy}\bigl(\log w(x)+\log w(y)\bigr).
\]
The right end is a strictly concave function with respect to \((c_{xy})_{x<y}\), and the constraint is linear, so the maximizer is unique.
Find the stationary point of the Lagrangian function and get for each \(x<y\),
\[
-2(\log c_{xy}+1)+\log w(x)+\log w(y)-2\lambda=0,
\]
Thus there exists a constant \(\alpha>0\) such that
\(
c_{xy}=\alpha\sqrt{w(x)w(y)}
\)。
From \(\sum_{x\ne y}c_{xy}=1\) and the identity
\(
\sum_{x<y}\sqrt{w(x)w(y)}=(A^2-1)/2=C/2
\)
We can get \(\alpha=1/C\), so
\(
Q_{xy}=c_{xy}/w(x)=\frac{1}{C}\sqrt{\frac{w(y)}{w(x)}}
\)
and is consistent with \(Q_w\) which corolles \ref{cor:pom-diagonal-rate-critical-continuous-time-generator}. Certification completed.
\end{proof}

\begin{corollary}[Closed and Auditable Cheeger Upper Bound]\label{cor:pom-diagonal-rate-critical-cheeger-auditable-upper-bound}
Following the corollary \ref{cor:pom-diagonal-rate-critical-continuous-time-generator}, let \(c_{xy}:=w(x)(Q_w)_{xy}\).
Define for any \(\varnothing\ne S\subsetneq X\)
\[
\mathrm{vol}(S):=\sum_{x\in S}w(x),\qquad
\mathrm{cut}(S):=\sum_{x\in S,\,y\notin S}c_{xy}.
\]
then a closed form is established
\[
\mathrm{cut}(S)=\frac{1}{C}\Bigl(\sum_{x\in S}\sqrt{w(x)}\Bigr)\Bigl(\sum_{y\notin S}\sqrt{w(y)}\Bigr).
\]
Let the (equivalent) Cheeger constant be defined as
\[
\Phi(w):=\min_{\varnothing\ne S\subsetneq X}\frac{\mathrm{cut}(S)}{\min\{\mathrm{vol}(S),\mathrm{vol}(S^c)\}}.
\]
Take any ordering \(w_{(1)}\ge\cdots\ge w_{(n)}\) and let \(S_k:=\{(1),\dots,(k)\}\).
Define sweep certificate
\[
\Phi_{\downarrow}(w):=\min_{1\le k\le n-1}\frac{\mathrm{cut}(S_k)}{\min\{\mathrm{vol}(S_k),\mathrm{vol}(S_k^c)\}},
\]
Then there is \(\Phi(w)\le \Phi_{\downarrow}(w)\).
Let \(\mu_2\) be the second eigenvalue of \(L_w\) (theorem \ref{thm:pom-diagonal-rate-critical-line-laplacian-first-order}), then it is obtained from the standard Cheeger upper bound of the reversible chain
\[
\mu_2\le 2\Phi(w)\le 2\Phi_{\downarrow}(w),
\]
where \(\Phi_{\downarrow}(w)\) can be computed explicitly in \(O(n)\) time after sorting \(w\).
\end{corollary}
\begin{proof}
cut closed form by
\(
c_{xy}=\frac{1}{C}\sqrt{w(x)w(y)}
\)
Obtained by direct summation.
By definition, \(\Phi(w)\) is the minimum value for all non-trivial subsets, so limiting the candidate set to \(\{S_k\}_{k=1}^{n-1}\) can only increase the minimum value, so \(\Phi(w)\le \Phi_{\downarrow}(w)\).
The last formula is the upper bound direction of the classic Cheeger inequality (for the spectral gap of the reversible continuous time chain), which is obtained by combining \(\Phi(w)\le \Phi_{\downarrow}(w)\).
The calculation of \(\Phi_{\downarrow}(w)\) requires only one sort and a linear scan of the prefix sum. Certification completed.
\end{proof}



\begin{theorem}[Secular equations of critical non-zero spectra interleaved with ordinal statistics]\label{thm:pom-diagonal-rate-critical-spectrum-secular-interlacing}
Follow the theorem \ref{thm:pom-diagonal-rate-critical-line-laplacian-first-order} The sign of \(n:=|X|\ge 2\), and note
\[
t_x:=\frac{A}{\sqrt{w(x)}}\qquad(x\in X).
\]
set up \(\tau_1<\cdots<\tau_m\) for \(\{t_x:x\in X\}\) are mutually distinct values, and \(\tau_j\) has a multiplicity of \(m_j\)（\(\sum_j m_j=n\)）。
but \(CL_w=\mathrm{diag}(t)-\mathbf 1\mathbf 1^{\mathsf T}\) consists of two types of parts:
\begin{enumerate}
  \item for each \(j\), eigenvalue \(\tau_j\) in multiplicity \(m_j-1\) Appear;
  \item rest \(m\) eigenvalues ​​are secular equations
  \[
  1-\sum_{x\in X}\frac{1}{t_x-\lambda}=0
  \]
  all real roots of \(\lambda_1=0\), and in each open interval \((\tau_j,\tau_{j+1})\) has exactly one root.
\end{enumerate}
Will \(\lambda=\mu C\), we get \(L_w\) positioning of non-zero eigenvalues; when \(t_x\) are different in pairs (i.e. \(m=n\)), the strictly interleaved relationship becomes
\[
\mu_{k+1}\in\Bigl(\frac{t_{(k)}}{C},\frac{t_{(k+1)}}{C}\Bigr)\qquad(k=1,\dots,n-1),
\]
in \(t_{(1)}<\cdots<t_{(n)}\) for \(t_x\) in ascending order.
\end{theorem}
\begin{proof}
Depend on \(CL_w=\mathrm{diag}(t)-\mathbf 1\mathbf 1^{\mathsf T}\) and its symmetry, the problem can be transformed into a rank-one perturbation spectrum analysis of a diagonal matrix.
\par
like \(\tau\) is the value of a certain diagonal element and its multiplicity is \(r\ge 2\), then any non-zero vector supported on the corresponding coordinate set and whose coordinate sum is zero \(v\),have \(\mathbf 1^{\mathsf T}v=0\) therefore \((\mathrm{diag}(t)-\mathbf 1\mathbf 1^{\mathsf T})v=\tau v\)。
According to the dimension counting, the dimension of the feature space is \(r-1\), so the first type of eigenvalues ​​is obtained.
\par
On the other hand, for \(\lambda\notin\{\tau_1,\dots,\tau_m\}\) using the rank-determinant lemma, we get
\[
0=\det\bigl(\mathrm{diag}(t)-\mathbf 1\mathbf 1^{\mathsf T}-\lambda I\bigr)
=\Bigl(\prod_{x\in X}(t_x-\lambda)\Bigr)\Bigl(1-\sum_{x\in X}\frac{1}{t_x-\lambda}\Bigr).
\]
because \(\lambda\notin\{t_x\}\), so the root must satisfy the secular equation.
remember \(F(\lambda):=1-\sum_x\frac{1}{t_x-\lambda}\),but \(F\) in each connected interval \((-\infty,\tau_1)\) and \((\tau_j,\tau_{j+1})\) continuously and strictly decreasing;
and when \(\lambda\uparrow\tau_{j+1}\) hour \(F(\lambda)\to-\infty\),when \(\lambda\downarrow\tau_j\) hour \(F(\lambda)\to+\infty\)。
Therefore in every \((\tau_j,\tau_{j+1})\) has exactly a unique zero point; and because
\(
\sum_x 1/t_x=\sum_x\sqrt{w(x)}/A=1
\), it can be seen \(F(0)=0\),thereby \(\lambda_1=0\) for \((-\infty,\tau_1)\) is the only zero point within.
This gives the completeness of the second type of eigenvalues.
Finally will \(\lambda=\mu C\) Restore to get it right \(L_w\); when \(t_x\) are different in pairs, the above interval is \((t_{(k)},t_{(k+1)})\), thus obtaining a strictly interleaved form. Certification completed.
\end{proof}

Theorem \ref{thm:pom-diagonal-rate-critical-spectrum-secular-interlacing} provides the secular equation for the critical nonzero spectrum.
Under product distributions, this one-dimensional secular function satisfies an exact multiplicative Stieltjes recursion in the tensor-product sense.

\begin{proposition}[Multiplicative Stieltjes recursion under product structure]\label{prop:pom-diagonal-rate-critical-spectrum-stieltjes-tensor-recursion}
For any full-support distribution \(w\in\Delta(X)\), write
\[
A(w):=\sum_{x\in X}\sqrt{w(x)},\qquad
C(w):=A(w)^2-1,
\qquad
t_x(w):=\frac{A(w)}{\sqrt{w(x)}}.
\]
Define the secular functions
\[
\mathcal S_w(z):=\sum_{x\in X}\frac{1}{t_x(w)-z},
\qquad
\mathcal F_w(z):=1-\mathcal S_w(z),
\]
whose zeros characterize the nontrivial spectrum of \(CL_w\) in Theorem \ref{thm:pom-diagonal-rate-critical-spectrum-secular-interlacing}.
\par
Let \(v\in\Delta(Y)\) be full support and \(\omega:=w\otimes v\in\Delta(X\times Y)\).
Then for any \((x,y)\in X\times Y\),
\(
t_{(x,y)}(\omega)=t_x(w)\,t_y(v)
\),
and for any complex \(z\) (away from poles) the exact identities hold:
\[
\mathcal S_{w\otimes v}(z)
=\sum_{x\in X}\frac{1}{t_x(w)}\,\mathcal S_v\!\Bigl(\frac{z}{t_x(w)}\Bigr),
\qquad
\mathcal F_{w\otimes v}(z)
=1-\sum_{x\in X}\frac{1}{t_x(w)}\Bigl(1-\mathcal F_v\!\bigl(z/t_x(w)\bigr)\Bigr).
\]
\end{proposition}
\begin{proof}
From \(A(w\otimes v)=A(w)A(v)\), one gets
\[
t_{(x,y)}(\omega)
=\frac{A(w\otimes v)}{\sqrt{\omega(x,y)}}
=\frac{A(w)A(v)}{\sqrt{w(x)v(y)}}
=t_x(w)\,t_y(v).
\]
Therefore,
\[
\mathcal S_{w\otimes v}(z)
=\sum_{x\in X}\sum_{y\in Y}\frac{1}{t_x(w)t_y(v)-z}
=\sum_{x\in X}\frac{1}{t_x(w)}\sum_{y\in Y}\frac{1}{t_y(v)-z/t_x(w)}
=\sum_{x\in X}\frac{1}{t_x(w)}\,\mathcal S_v\!\Bigl(\frac{z}{t_x(w)}\Bigr).
\]
The identity for \(\mathcal F\) follows by substitution from \(\mathcal F=1-\mathcal S\). This completes the proof.
\end{proof}

\begin{corollary}[Iterative composition for \texorpdfstring{$k$}{k}-fold tensors]\label{cor:pom-diagonal-rate-critical-spectrum-stieltjes-tensor-recursion-kfold}
Take \(k\ge 2\) full-support distributions \(w^{(1)}\in\Delta(X_1),\dots,w^{(k)}\in\Delta(X_k)\), and let \(\omega:=\bigotimes_{i=1}^k w^{(i)}\).
Then for any complex \(z\) (away from poles), one has the iterative expansion
\[
\mathcal S_{\omega}(z)
=\sum_{x_1\in X_1}\frac{1}{t_{x_1}(w^{(1)})}
\sum_{x_2\in X_2}\frac{1}{t_{x_2}(w^{(2)})}\cdots
\sum_{x_{k-1}\in X_{k-1}}\frac{1}{t_{x_{k-1}}(w^{(k-1)})}\,
\mathcal S_{w^{(k)}}\!\Bigl(\frac{z}{t_{x_1}(w^{(1)})\cdots t_{x_{k-1}}(w^{(k-1)})}\Bigr).
\]
\end{corollary}
\begin{proof}
Apply Proposition \ref{prop:pom-diagonal-rate-critical-spectrum-stieltjes-tensor-recursion} repeatedly and expand recursively. This completes the proof.
\end{proof}



\begin{corollary}[Closed form of binary and uniform distribution]\label{cor:pom-diagonal-rate-critical-spectrum-binary-uniform}
Follow the theorem \ref{thm:pom-diagonal-rate-critical-line-laplacian-first-order} of \(L_w\)。
\begin{enumerate}
  \item if \(n=2\) and \(w=(p,1-p)\), then the only non-zero eigenvalue is
  \[
  \mu_2=\frac{1}{2p(1-p)}.
  \]
  \item if \(w\equiv 1/n\),but
  \[
  \mu_2=\cdots=\mu_n=\frac{n}{n-1}.
  \]
\end{enumerate}
So in the uniform case \(1-\lambda_2(K_\delta)=\frac{n}{n-1}\delta+o(\delta)\), in the binary case \(1-\lambda_2(K_\delta)=\frac{\delta}{2p(1-p)}+o(\delta)\)。
\end{corollary}
\begin{proof}
\textup{(1)} Depend on \(A=\sqrt p+\sqrt{1-p}\)、\(C=A^2-1=2\sqrt{p(1-p)}\) have to
\(
t_1=A/\sqrt p,\ t_2=A/\sqrt{1-p}
\)。
matrix \(CL_w=\mathrm{diag}(t)-\mathbf 1\mathbf 1^{\mathsf T}\) in the second-order case, the non-zero eigenvalue is equal to the trace
\[
\lambda_2=t_1+t_2-2
=\Bigl(\sqrt p+\sqrt{1-p}\Bigr)\Bigl(\frac{1}{\sqrt p}+\frac{1}{\sqrt{1-p}}\Bigr)-2
=\frac{1}{\sqrt{p(1-p)}},
\]
thereby \(\mu_2=\lambda_2/C=1/(2p(1-p))\)。
\par
\textup{(2)} like \(w\equiv 1/n\),but \(A=\sqrt n\)、\(C=n-1\)、\(t_x=n\) constant, so \(CL_w=nI-\mathbf 1\mathbf 1^{\mathsf T}\)。
Its in \(\mathbf 1^\perp\) above is \(nI\),exist \(\mathrm{span}\{\mathbf 1\}\) is zero, so \(L_w\) The non-zero spectrum is \(n/(n-1)\) (multiplicity \(n-1\)）。
The first order law of the final spectral gap is given by
\(
\lambda_2(K_\delta)=1-\delta\mu_2+o(\delta)
\)
Launch immediately. Certification completed.
\end{proof}

\begin{theorem}[Closed form of pseudo-determinant of critical line Laplace]\label{thm:pom-diagonal-rate-critical-laplacian-pseudodeterminant}
Follow the theorem \ref{thm:pom-diagonal-rate-critical-line-laplacian-first-order} And order \(n:=|X|\)。
Define pseudo-determinant (non-zero spectral product)
\[
\operatorname{pdet}(L_w):=\prod_{i=2}^{n}\mu_i.
\]
Then there is a strict identity
\[
\operatorname{pdet}(L_w)=\frac{A^{\,n-2}}{C^{\,n-1}}\cdot \frac{1}{\sqrt{\prod_{x\in X}w(x)}}.
\]
Equivalently, note the geometric mean
\(
G(w):=\exp\bigl(\frac1n\sum_x\log w(x)\bigr)
\)
and \(H_{1/2}(w)=2\log A\)(proposition \ref{prop:pom-renyi-half-hellinger-identity-tensor-additivity}),but
\[
\log\operatorname{pdet}(L_w)
=\frac{n-2}{2}H_{1/2}(w)
-(n-1)\log\bigl(e^{H_{1/2}(w)}-1\bigr)
-\frac{n}{2}\log G(w).
\]
\end{theorem}
\begin{proof}
remember \(L':=CL_w=\mathrm{diag}(t)-\mathbf 1\mathbf 1^{\mathsf T}\)。
right \(\lambda>0\) using the rank-determinant lemma, we get
\[
\det(L'+\lambda I)=\Bigl(\prod_{x\in X}(t_x+\lambda)\Bigr)\Bigl(1-\sum_{x\in X}\frac{1}{t_x+\lambda}\Bigr).
\]
Depend on \(\sum_x 1/t_x=1\) writable
\[
1-\sum_{x}\frac{1}{t_x+\lambda}
=\sum_x\Bigl(\frac{1}{t_x}-\frac{1}{t_x+\lambda}\Bigr)
=\lambda\sum_x\frac{1}{t_x(t_x+\lambda)}.
\]
then
\[
\lim_{\lambda\downarrow 0}\frac{\det(L'+\lambda I)}{\lambda}
=\Bigl(\prod_x t_x\Bigr)\sum_x\frac{1}{t_x^2}.
\]
on the other hand,\(\operatorname{pdet}(L')=\prod_{i=2}^n(C\mu_i)=C^{n-1}\operatorname{pdet}(L_w)\),and
\(
\operatorname{pdet}(L')=\lim_{\lambda\downarrow 0}\det(L'+\lambda I)/\lambda
\)。
therefore
\[
\operatorname{pdet}(L_w)
=\frac{1}{C^{n-1}}\Bigl(\prod_x t_x\Bigr)\sum_x\frac{1}{t_x^2}.
\]
Will \(t_x=A/\sqrt{w(x)}\) Substitute,
\[
\prod_x t_x=\frac{A^{\,n}}{\sqrt{\prod_x w(x)}},
\qquad
\sum_x\frac{1}{t_x^2}=\sum_x\frac{w(x)}{A^2}=\frac{1}{A^2},
\]
thus getting \(\operatorname{pdet}(L_w)=A^{n-2}/(C^{n-1}\sqrt{\prod w})\)。
The logarithmic form is given by \(H_{1/2}(w)=2\log A\)、\(C=e^{H_{1/2}(w)}-1\) and \(\sqrt{\prod w}=G(w)^{n/2}\) is a direct simplification. Certification completed.
\end{proof}

Theorem \ref{thm:pom-diagonal-rate-critical-laplacian-pseudodeterminant} gives a closed form for \(\operatorname{pdet}(L_w)\).
Under tensor products, this closed form induces a fully explicit synergy-factor decomposition.

\begin{theorem}[Synergy-factor decomposition of the pseudodeterminant under tensor products]\label{thm:pom-diagonal-rate-critical-laplacian-pdet-tensor-synergy-factorization}
Take two finite sets \(X,Y\), full-support distributions \(w\in\Delta(X)\), \(v\in\Delta(Y)\), and let \(\omega:=w\otimes v\in\Delta(X\times Y)\).
Write \(n:=|X|\ge 2\), \(m:=|Y|\ge 2\).
Define
\[
A(w):=\sum_{x\in X}\sqrt{w(x)},\quad C(w):=A(w)^2-1,
\qquad
A(v):=\sum_{y\in Y}\sqrt{v(y)},\quad C(v):=A(v)^2-1.
\]
Then the exact identity holds:
\[
\operatorname{pdet}(L_{\omega})
=\operatorname{pdet}(L_w)^{\,m}\operatorname{pdet}(L_v)^{\,n}\cdot \mathsf S(w,v),
\]
where the synergy factor \(\mathsf S(w,v)\) depends only on \(n,m\) and \(C(w),C(v)\), with closed form
\[
\mathsf S(w,v)
:=\frac{(1+C(w))^{m-1}(1+C(v))^{n-1}\,C(w)^{m(n-1)}\,C(v)^{n(m-1)}}
{\bigl(C(w)+C(v)+C(w)C(v)\bigr)^{nm-1}}.
\]
\end{theorem}
\begin{proof}
By Theorem \ref{thm:pom-diagonal-rate-critical-laplacian-pseudodeterminant},
\[
\operatorname{pdet}(L_w)=\frac{A(w)^{n-2}}{C(w)^{n-1}}\cdot \frac{1}{\sqrt{\prod_{x\in X}w(x)}},
\qquad
\operatorname{pdet}(L_v)=\frac{A(v)^{m-2}}{C(v)^{m-1}}\cdot \frac{1}{\sqrt{\prod_{y\in Y}v(y)}}.
\]
Also,
\[
A(\omega)=A(w)A(v),\qquad
C(\omega)=A(\omega)^2-1=C(w)+C(v)+C(w)C(v),
\]
and
\[
\prod_{(x,y)\in X\times Y}\omega(x,y)
=\Bigl(\prod_{x\in X}w(x)\Bigr)^{m}\Bigl(\prod_{y\in Y}v(y)\Bigr)^{n}.
\]
Substituting these identities into \(\operatorname{pdet}(L_\omega)\) and \(\operatorname{pdet}(L_w)^m\operatorname{pdet}(L_v)^n\), canceling geometric-mean terms, and collecting exponents gives
\[
\frac{\operatorname{pdet}(L_{\omega})}{\operatorname{pdet}(L_w)^{\,m}\operatorname{pdet}(L_v)^{\,n}}
=\frac{A(w)^{2(m-1)}A(v)^{2(n-1)}\,C(w)^{m(n-1)}\,C(v)^{n(m-1)}}
{C(\omega)^{nm-1}}.
\]
Using \(A(w)^2=1+C(w)\) and \(A(v)^2=1+C(v)\) yields the stated closed form for the synergy factor. This completes the proof.
\end{proof}

\begin{corollary}[Pure \texorpdfstring{$H_{1/2}$}{H1/2} structure after geometric-mean stripping]\label{cor:pom-diagonal-rate-critical-laplacian-pdet-tensor-synergy-geometricmean-stripped}
Using the geometric mean from Theorem \ref{thm:pom-diagonal-rate-critical-laplacian-pseudodeterminant},
\(
G(w):=\exp(\frac1n\sum_{x\in X}\log w(x))
\),
define the rescaled pseudodeterminant
\[
\widetilde{\operatorname{pdet}}(L_w):=\operatorname{pdet}(L_w)\,G(w)^{n/2}.
\]
Then for the product distribution \(\omega=w\otimes v\),
\[
\widetilde{\operatorname{pdet}}(L_{w\otimes v})
=\widetilde{\operatorname{pdet}}(L_w)^{\,m}\widetilde{\operatorname{pdet}}(L_v)^{\,n}\cdot \mathsf S(w,v),
\]
where \(\mathsf S(w,v)\) is the closed form from Theorem \ref{thm:pom-diagonal-rate-critical-laplacian-pdet-tensor-synergy-factorization}, and via
\(
C(\cdot)=e^{H_{1/2}(\cdot)}-1
\)
it can be expressed solely in terms of \(H_{1/2}\).
\end{corollary}
\begin{proof}
From \(\prod_{(x,y)}\omega(x,y)=(\prod_x w(x))^m(\prod_y v(y))^n\), we immediately get \(G(w\otimes v)=G(w)\,G(v)\).
Therefore, multiplying both sides of Theorem \ref{thm:pom-diagonal-rate-critical-laplacian-pdet-tensor-synergy-factorization} by \(G(w\otimes v)^{nm/2}=G(w)^{mn/2}G(v)^{nm/2}\) gives the stated rescaled identity.
The final statement is the direct rewrite \(C=e^{H_{1/2}}-1\). This completes the proof.
\end{proof}



\begin{corollary}[Natural constants of the uniform limit]\label{cor:pom-diagonal-rate-critical-laplacian-pdet-uniform-limit-e}
like \(w\equiv 1/n\),but
\[
\operatorname{pdet}(L_w)=\Bigl(\frac{n}{n-1}\Bigr)^{n-1}\xrightarrow[n\to\infty]{}e.
\]
\end{corollary}
\begin{proof}
By theorem \ref{thm:pom-diagonal-rate-critical-laplacian-pseudodeterminant} Substitute \(A=\sqrt n\)、\(C=n-1\)、\(\prod_x w(x)=n^{-n}\) directly obtains the closed form; the limit is given by \((1+1/(n-1))^{n-1}\to e\) is given. Certification completed.
\end{proof}

\begin{theorem}[Explicit closure and scaling of critical polynomials \texorpdfstring{$\zeta$}{zeta} limit]\label{thm:pom-diagonal-rate-critical-xi-polynomial-and-zeta-scaling}
Follow the theorem \ref{thm:pom-diagonal-rate-critical-line-laplacian-first-order},make
\[
\Xi_w(s):=\det{}'(L_w+sI)=\prod_{i=2}^{n}(\mu_i+s),\qquad s\in\RR,
\]
in \(\det{}'\) means removing zero eigenvalue factors (equivalent to \(\det(L_w+sI)/s\)）。
Then for any \(s\in\CC\) has a closed form
\[
\Xi_w(s)
=\frac{1}{A\,C^{\,n-1}}\Bigl(\prod_{x\in X}\frac{1}{\sqrt{w(x)}}\Bigr)
\sum_{x\in X} w(x)\prod_{y\ne x}\bigl(A+sC\sqrt{w(y)}\bigr),
\]
thereby \(\Xi_w\) is the number of times \(n-1\) is a polynomial of one variable and satisfies \(\Xi_w(0)=\operatorname{pdet}(L_w)\) (and theorem \ref{thm:pom-diagonal-rate-critical-laplacian-pseudodeterminant} consistent).
\par
Furthermore, remember that the finite dimension \(\zeta\) determinant
\(
\zeta_\delta(z):=\det(I-zK_\delta)^{-1}
\)
(Same corollary \ref{cor:pom-diagonal-rate-rank-one-refresh-trace-zeta-fingerprint}）。
Then for any fixed \(s\in\CC\) has a zoom limit
\[
\lim_{\delta\downarrow 0}\ \delta^{-(n-1)}\,
\frac{\det(I-e^{-s\delta}K_\delta)}{1-e^{-s\delta}}
=\Xi_w(s).
\]
\end{theorem}
\begin{proof}
Depend on \(L'=CL_w=\mathrm{diag}(t)-\mathbf 1\mathbf 1^{\mathsf T}\) and the rank determinant lemma,
\[
\det(L'+sCI)
=\Bigl(\prod_x(t_x+sC)\Bigr)\Bigl(1-\sum_x\frac{1}{t_x+sC}\Bigr).
\]
Notice \(\sum_x 1/t_x=1\), so
\[
1-\sum_x\frac{1}{t_x+sC}
=\sum_x\Bigl(\frac{1}{t_x}-\frac{1}{t_x+sC}\Bigr)
=sC\sum_x\frac{1}{t_x(t_x+sC)}.
\]
Substitute it and use \(L_w=L'/C\) have to
\[
\Xi_w(s)=\frac{\det(L_w+sI)}{s}
=\frac{1}{C^{n}}\cdot \frac{1}{s}\,\det(L'+sCI)
=\frac{1}{C^{n-1}}\Bigl(\prod_x(t_x+sC)\Bigr)\sum_x\frac{1}{t_x(t_x+sC)}.
\]
Substitute \(t_x=A/\sqrt{w(x)}\) and dividing the denominators together, we get
\[
\Xi_w(s)=\frac{1}{A\,C^{n-1}}\Bigl(\prod_{x}\frac{1}{\sqrt{w(x)}}\Bigr)\sum_{x}w(x)\prod_{y\ne x}\bigl(A+sC\sqrt{w(y)}\bigr),
\]
It is thus a polynomial of one variable and the constant term is given by \(s=0\) and theorem \ref{thm:pom-diagonal-rate-critical-laplacian-pseudodeterminant} Consistency is given immediately.
\par
The scaling limit is given by similarity invariance
\(
S_\delta=I-\delta L_w+o(\delta)
\)
(theorem \ref{thm:pom-diagonal-rate-critical-line-laplacian-first-order})and
\(
e^{-s\delta}=1-s\delta+O(\delta^2)
\)
roll out:
\[
I-e^{-s\delta}S_\delta
=I-\bigl(1-s\delta+O(\delta^2)\bigr)\bigl(I-\delta L_w+o(\delta)\bigr)
=\delta(L_w+sI)+o(\delta),
\]
thereby
\(
\det(I-e^{-s\delta}K_\delta)=\det(I-e^{-s\delta}S_\delta)=\delta^{n}\det(L_w+sI)+o(\delta^{n})
\)
and
\(
1-e^{-s\delta}=s\delta+O(\delta^2)
\)。
Divide the two and multiply by \(\delta^{-(n-1)}\) get the limit \(\Xi_w(s)=\det(L_w+sI)/s\). Certification completed.
\end{proof}

\begin{corollary}[Interleaving of critical spectra and clamping of the first spectral gap]\label{cor:pom-diagonal-rate-critical-spectrum-gap-bounds}
Follow the theorem \ref{thm:pom-diagonal-rate-critical-spectrum-secular-interlacing} And order \(w_{\max}:=\max_x w(x)\)，\(w_{\mathrm{2nd}}\) is the second largest mass (taken when there is a parallel maximum \(w_{\mathrm{2nd}}=w_{\max}\)）。
Then the minimum non-zero eigenvalue satisfies
\[
\frac{A}{C\sqrt{w_{\max}}}\ \le\ \mu_2\ \le\ \frac{A}{C\sqrt{w_{\mathrm{2nd}}}},
\]
And when the maximum mass is unique, both ends are strict inequalities.
\end{corollary}
\begin{proof}
By theorem \ref{thm:pom-diagonal-rate-critical-spectrum-secular-interlacing} staggered positioning,\(\mu_2 C\) is located between the smallest two levels of mutually different diagonal values.
and \(t_x=A/\sqrt{w(x)}\) follow \(w(x)\) strictly decreases, so the minimum diagonal value corresponds to \(w_{\max}\), the second smallest diagonal value corresponds to \(w_{\mathrm{2nd}}\)。
Divide the endpoints of the interval by \(C\) is obtained. Certification completed.
\end{proof}

\begin{proposition}[Explicit shape and sign separation of critical eigenvectors]\label{prop:pom-diagonal-rate-critical-eigenvector-shape-sign-separation}
set up \(t_x\) are different in pairs, and let \(\mu_{k+1}\) is the theorem \ref{thm:pom-diagonal-rate-critical-spectrum-secular-interlacing} middle range \((t_{(k)}/C,t_{(k+1)}/C)\) corresponding non-zero eigenvalues.
remember \(\lambda_{k+1}:=\mu_{k+1}C\in(t_{(k)},t_{(k+1)})\), and order \(v^{(k+1)}\) is the matrix \(L':=CL_w=\mathrm{diag}(t)-\mathbf 1\mathbf 1^{\mathsf T}\) corresponding eigenvector (\(\lambda_{k+1}\)-eigenvector).
then there is a constant \(\sigma\ne 0\) applies to all \(x\in X\) have
\[
v^{(k+1)}_x=\frac{\sigma}{t_x-\lambda_{k+1}}.
\]
and pressing \(t_{(1)}<\cdots<t_{(n)}\) in the sorted coordinates,\(\bigl(v^{(k+1)}_{(i)}\bigr)_{i=1}^n\) symbol is in \(k\) and \(k+1\) occurs uniquely between: that is, the components are in \(i\le k\) and \(i\ge k+1\) has the same sign on both sides and opposite signs (the overall sign can be \(\sigma\) choose).
\end{proposition}
\begin{proof}
Depend on \((\mathrm{diag}(t)-\mathbf 1\mathbf 1^{\mathsf T})v=\lambda v\) have to
\[
(t_x-\lambda)v_x=\sum_{y\in X}v_y=:S,
\]
right end and \(x\) has nothing to do, so \(v_x=S/(t_x-\lambda)\)。
like \(S=0\),but \((t_x-\lambda)v_x=0\) for all \(x\) is established because \(t_x\) are pairwise different and \(\lambda\notin\{t_x\}\) (by theorem \ref{thm:pom-diagonal-rate-critical-spectrum-secular-interlacing}),force \(v=0\),contradiction.
Therefore \(S\ne 0\),Pick \(\sigma:=S\) to get an explicit shape.
\par
Separation of symbols, because \(\lambda_{k+1}\in(t_{(k)},t_{(k+1)})\),have \(t_{(i)}-\lambda_{k+1}<0\) when \(i\le k\),and \(t_{(i)}-\lambda_{k+1}>0\) when \(i\ge k+1\)。
Therefore \(v_{(i)}\propto (t_{(i)}-\lambda_{k+1})^{-1}\) has opposite signs on both sides, and there are no additional flips other than the overall sign. Certification completed.
\end{proof}

\begin{proposition}[Closed expression of critical spectral invariant]\label{prop:pom-diagonal-rate-critical-spectrum-invariants-closed}
Follow the theorem \ref{thm:pom-diagonal-rate-critical-line-laplacian-first-order} And order \(n:=|X|\)。
Define the sum of inverse square roots
\[
B(w):=\sum_{x\in X}w(x)^{-1/2},
\qquad
p_{-1}(w):=\sum_{x\in X}w(x)^{-1}.
\]
Then the critical spectrum \(\{\mu_i\}_{i=2}^n\) satisfies the following strict identities:
\[
\sum_{i=2}^{n}\mu_i=\tr(L_w)=\frac{A\,B(w)-n}{C},
\qquad
\prod_{i=2}^{n}\mu_i=\operatorname{pdet}(L_w)=\frac{A^{\,n-2}}{C^{\,n-1}\sqrt{\prod_x w(x)}},
\]
And the sum of quadratic powers satisfies
\[
\sum_{i=2}^{n}\mu_i^2=\tr(L_w^2)=\frac{A^2\,p_{-1}(w)-2A\,B(w)+n^2}{C^2}.
\]
Therefore, the polynomial \(\Xi_w(s)=\prod_{i=2}^{n}(\mu_i+s)\) are simultaneously encoded with low-order coefficients \(A\) (equivalent to \(H_{1/2}(w)\)), geometric mean \(\prod_x w(x)\) as well as \(B(w),p_{-1}(w)\) and other inverse moment information.
\end{proposition}
\begin{proof}
Depend on \(L_w=(A\,\mathrm{diag}(w^{-1/2})-\mathbf 1\mathbf 1^{\mathsf T})/C\), whose trace is
\[
\tr(L_w)=\frac{1}{C}\Bigl(A\sum_x w(x)^{-1/2}-\tr(\mathbf 1\mathbf 1^{\mathsf T})\Bigr)
=\frac{A\,B(w)-n}{C}.
\]
The pseudo-determinant identity is given by the theorem \ref{thm:pom-diagonal-rate-critical-laplacian-pseudodeterminant} given.
\par
for calculation \(\tr(L_w^2)\),remember \(D:=\mathrm{diag}(w^{-1/2})\)、\(J:=\mathbf 1\mathbf 1^{\mathsf T}\),but
\(
L_w=\frac{1}{C}(AD-J)
\)。
Expand and trace
\[
\tr(L_w^2)
=\frac{1}{C^2}\tr\bigl(A^2D^2-A(DJ+JD)+J^2\bigr).
\]
Notice \(\tr(D^2)=\sum_x w(x)^{-1}=p_{-1}(w)\),and \(\tr(DJ)=\tr(JD)=\sum_x w(x)^{-1/2}=B(w)\),and \(J^2=nJ\) thus \(\tr(J^2)=n\tr(J)=n^2\)。
Substituting back we get the quadratic powers and identities shown. Certification completed.
\end{proof}

